% Set up various types of listings to be used in different files.
% This provides a uniform style for listings in all documents

% define a listing postbreak character to be used
% Note that this version works safely with cmyk mode for xcolors.
% It uses the STIX Match font to have the correct character at the LuaLaTeX engine level.
\newsavebox{\lstbreakbox}
\sbox{\lstbreakbox}{\textcolor{red}{\mystixmathfont ↪}\space}

\lstdefinelanguage{BibTeX}{
 basicstyle=\footnotesize\ttfamily,
  keywordstyle=\color{blue}\bfseries,      % @InProceedings in Blue
  keywordstyle=[2]\color{darkgray},       % author, title, etc. in Gray
  stringstyle=\color{red},                 % "Values" in Red
  commentstyle=\color{green!50!black},
  showstringspaces=false,
  breaklines=true,
  postbreak=\usebox{\lstbreakbox},
  keywords={%
    @article,@book,@collectedbook,@conference,@electronic,@ieeetranbstctl,%
    @inbook,@incollectedbook,@incollection,@injournal,@inproceedings,%
    @manual,@mastersthesis,@misc,@patent,@periodical,@phdthesis,@preamble,%
    @proceedings,@standard,@string,@techreport,@unpublished%
  },
  % Entry fields/keys (e.g., author, title)
  keywords=[2]{%
    address,author,booktitle,chapter,edition,editor,howpublished,institution,%
    journal,key,month,note,number,organization,pages,publisher,school,%
    series,title,type,volume,year%
  },
  comment=[l][\itshape]{@comment},
  sensitive=false,
  % String handling
  morestring=[b]",      % Matches "double quotes"
  % morestring=[s]{\{}{\}}  % The correct way to handle {balanced braces}
  % FIX: use moredelim with a style argument and escaped braces
  % moredelim=[s][\color{red}]{\{}{\}}
}

\lstdefinestyle{latexExampleForAuthors}{
language=[LaTeX]{TeX},
    breaklines=true,
    postbreak=\usebox{\lstbreakbox},
    basicstyle=\small\tt,
    keywordstyle=\color{blue}\sf,
    identifierstyle=\color{magenta},
    commentstyle=\color{cyan},
    backgroundcolor=\color{yellow!15},
    extendedchars=true,
    inputencoding=utf8,
    tabsize=2,
    columns=flexible,
    morekeywords={subtitle, alttitle, altsubtitle, hostcompany, courseCycle, courseCode, courseCredits, programcode, degreeName, subjectArea, educationSubjectcode, nationalsubjectcategories, todo, subsection}
}


\lstdefinestyle{latexExample}{
language=[LaTeX]{TeX},
    breaklines=true,
    postbreak=\usebox{\lstbreakbox},
    basicstyle=\small\tt,
    keywordstyle=\color{blue}\sf,
    identifierstyle=\color{magenta},
    commentstyle=\color{cyan},
    backgroundcolor=\color{yellow!15},
    tabsize=2,
    columns=flexible,
}

%
% The commands below are to configure JSON listings
% 
% format for JSON listings
\colorlet{punct}{red!60!black}
\definecolor{delim}{RGB}{20,105,176}
\definecolor{numb}{RGB}{106, 109, 32}
\definecolor{string}{RGB}{0, 0, 0}

\lstdefinelanguage{json}{
    numbers=none,
    numberstyle=\small,
    frame=none,
    rulecolor=\color{black},
    showspaces=false,
    showtabs=false,
    breaklines=true,
    postbreak=\usebox{\lstbreakbox},
    breakatwhitespace=true,
    basicstyle=\ttfamily\small,
    extendedchars=false,
    upquote=true,
    morestring=[b]",
    stringstyle=\color{string},
    literate=
     *{0}{{{\color{numb}0}}}{1}
      {1}{{{\color{numb}1}}}{1}
      {2}{{{\color{numb}2}}}{1}
      {3}{{{\color{numb}3}}}{1}
      {4}{{{\color{numb}4}}}{1}
      {5}{{{\color{numb}5}}}{1}
      {6}{{{\color{numb}6}}}{1}
      {7}{{{\color{numb}7}}}{1}
      {8}{{{\color{numb}8}}}{1}
      {9}{{{\color{numb}9}}}{1}
      {:}{{{\color{punct}{:}}}}{1}
      {,}{{{\color{punct}{,}}}}{1}
      {\{}{{{\color{delim}{\{}}}}{1}
      {\}}{{{\color{delim}{\}}}}}{1}
      {\[}{{{\color{delim}{\[}}}}{1}
      {\]}{{{\color{delim}{\]}}}}{1}
      {’}{{\char13}}1,
}

\lstdefinelanguage{XML}
{
  basicstyle=\ttfamily\color{blue}\bfseries\small,
  morestring=[b]",
  morestring=[s]{>}{<},
  morecomment=[s]{<?}{?>},
  stringstyle=\color{black},
  identifierstyle=\color{blue},
  keywordstyle=\color{cyan},
  breaklines=true,
  postbreak=\usebox{\lstbreakbox},
  breakatwhitespace=true,
  morekeywords={xmlns,xsi,noNamespaceSchemaLocation,version,type,id,x,y,source,target,tool,transRef,roleRef,objective,eventually}% list your attributes here
}
\lstset{style=latexExample}

\lstdefinestyle{PDF}{
language=[LaTeX]{TeX},
    breaklines=true,
    postbreak=\usebox{\lstbreakbox},
    basicstyle=\small\tt,
    keywordstyle=\color{blue}\sf,
    identifierstyle=\color{magenta},
    commentstyle=\color{cyan},
    backgroundcolor=\color{yellow!15},
    extendedchars=true,
    inputencoding=utf8,
    stringstyle=\ttfamily,
    morestring=[b]',
    tabsize=2,
    columns=flexible,
    morekeywords={Type, Action, JavaScript, obj, endobj, Length, Filter, FlateDecode, stream, endstream, JS, ASCIIHexDecode, startxref, xref, trailer, Size, Root, Catalog, OpenAction, AcroForm, Pages, Outlines, S,
    Kids, Count, Page, Parent, MediaBox, Contents, Resources, Font, ProcSet, PDF, Text, F1}
}

% Extend the XML style
\lstdefinestyle{myXML}{
    language=XML,
    columns=flexible,
    breaklines=true,
    showspaces=false,                       % Dont make spaces visible
    showtabs=false,                         % Dont make tabs visible
    breakatwhitespace=true,
    postbreak=\usebox{\lstbreakbox},
    basicstyle=\ttfamily\small,
    columns=fixed,
    escapechar = £,
    morekeywords={dc:format, dc:title, dc:date, dc:type, dc:creator, dc:description, dc:subject, dc:source, dc:language, rdf:Alt, rdf:li, rdf:Bag, rdf:Seq, rdf:Description, rdf:Description, rdf:RDF, pdfaExtension:schemas, pdfaSchema:property, pdfaSchema:valueType, pdf:Producer, pdf:Keywords, pdf:PDFVersion, pdfaSchema:schema, pdfaSchema:prefix, pdfaSchema:namespaceURI, pdfaProperty:name, pdfaProperty:valueType, pdfaProperty:category, pdfaProperty:description,  xmp:CreateDate, xmp:ModifyDate, xmp:MetadataDate, xmp:CreatorTool, xmpMM:DocumentID, xmpMM:InstanceID, xmpMM:VersionID, xmpMM:RenditionClass, Iptc4xmpCore:CreatorContactInfo, Iptc4xmpCore:CiEmailWork, prism:complianceProfile, prism:pageCount, xmpTPg:NPages,  x:xmpmeta},
    %string=[bd]{|}      % Uncommenting this line eliminates the visible spaces in strings
% Support for Swedish, German and Portuguese umlauts
  literate=%
  {Ö}{{\"O}}1
  {Ä}{{\"A}}1
  {Å}{{\AA{}}}1
  {Ü}{{\"U}}1
  {ß}{{\ss}}1
  {ü}{{\"u}}1
  {ö}{{\"o}}1
  {ä}{{\"a}}1
  {å}{{\aa{}}}1
  {á}{{\'a}}1
  {ã}{{\~a}}1
  {é}{{\'e}}1
  {è}{{\`e}}1
  {€}{\euro}1%
  {’}{{\char13}}1
  {-}{{\textendash}}1
  {–}{{\textendash}}1
  {…}{{\ldots}}1,
  %{}{{BOM}{\;}{0xfeff}}{10},
}
